\chapter{Stand der Technik}\label{ch:sota}
% [AP-S1] Instanz-Spiegel von Kapitel 2 (Autor-Direktive 2026-06-16): jede Grundlagen-Definitionsklasse
% (Kap. 2) erhaelt hier ihre wissenschaftlich nachweisbaren Instanzen. 3.1 + 3.2 sind ausformuliert; die
% darunter folgenden Cluster-Sektionen (Cluster A--F + 3 Tabellen + Forschungsluecke) sind ROHMATERIAL und
% werden in den naechsten AP-S1-Inkrementen in 3.3 (Achsen-Sezierung) / 3.4 (Workloads) / 3.5 (Messen) /
% 3.6 (Architektur-Design) reorganisiert; die Forschungsluecke wird zu 3.7. Bis dahin vorlaeufige Nummerierung.

Wo Kapitel~\ref{ch:fundamentals} die konzeptionellen \emph{Definitions-Klassen} bereitstellt, führt
dieser Stand der Technik ihre wissenschaftlich nachweisbaren \emph{Instanzen} an: zu jeder
Grundlagen-Klasse die konkreten Veröffentlichungen, Algorithmen und Technologien, die sie umsetzen.
Beide Kapitel greifen damit unmittelbar ineinander (Abschnitt~\ref{ssec:measurement-patterns}).

\section{Überblick und Sezierungs-Prinzip}\label{sec:sota-overview}
Die Analyse stützt sich auf 33 tiefenanalysierte Suchalgorithmus-Veröffentlichungen (P01--P33) und einen
parallelen Allokator-Korpus aus 23 Arbeiten (A01--A23); die Achsen-Belegung jeder Arbeit ist als
XML-Profil in der \texttt{comdare-cache-engine} (\texttt{algorithm\_profiles/}) persistiert, die
vollständigen Bausteine- und Allokator-Matrizen stehen in Anhang~\ref{app:blocks}. Die
Veröffentlichungen entstammen sechs thematischen \emph{Herkunfts-Clustern} --- Trie-Familie, Hybrid-
und B\textsuperscript{+}-Familie, Layout-Theorie, Prefetching, Synchronisation und den hardware-nahen
TU-Dresden-Arbeiten ---; diese Cluster dienen hier nur der Herkunfts-Einordnung, nicht als
Gliederungsachse.

Der Kernbeitrag dieser Arbeit ist die \emph{Sezierung} (Abschnitt~\ref{ssec:own-framework}): Jeder
Gesamt-Algorithmus (ein „Lebewesen") wird in seine orthogonalen Organe --- die Achsen --- zerlegt, und
die so gewonnenen Organ-Algorithmen werden \emph{je Achse} einsortiert, mit kurzem Verweis auf das
Herkunfts-Lebewesen. Diese achsenweise Sezierung bildet den Kern des Kapitels; die Zuordnung
Achse\,$\leftrightarrow$\,Algorithmus\,$\leftrightarrow$\,Veröffentlichung folgt einer web-verifizierten
Provenienz-Map über rund 110 Organ-Algorithmen.

Dass das Achsen-Gerüst auch \emph{nicht-baumartige} Verfahren trägt, belegt eine reine Hash-Tabelle
(\texttt{absl::flat\_hash\_map}/SwissTable~\cite{abseil_swisstable}) als Gegenprobe innerhalb der
SearchAlgorithm-Gattung: Sie erfüllt dasselbe \texttt{std::map}-Interface über ein \emph{limitiertes
Achsen-Subset} --- die trie-spezifischen Achsen (Pfadkompression, Knotentyp, Bereichs-Index) entfallen
als Durchreich-Variante, während Allokation, Prefetch, Nebenläufigkeit, Speicher-Layout und ISA-SIMD
(für das SwissTable-Gruppen-Lookup) geteilt bleiben --- und erfordert so keine eigene Gattung.

\section{Cache-Konzepte als Instanzen}\label{sec:sota-cache}
Die Cache-Hierarchie-Konzepte aus Abschnitt~\ref{sec:cache-basics} treten im Stand der Technik als
konkrete Entwürfe und Mechanismen auf; jeder der folgenden Unterabschnitte spiegelt eine
Grundlagen-Definition aus Abschnitt~\ref{sec:cache-basics}.

\subsection{Speicher- und Cache-Ebenen}\label{ssec:sota-levels}
Die mehrstufige Hierarchie (Abschnitt~\ref{ssec:cache-levels}) ist Gegenstand grundlegender
Vermessungen: Graefe und Larson (P15, 2001) untersuchen B-Baum-Indizes explizit gegenüber den
CPU-Cache-Ebenen, Bender, Demaine und Farach-Colton (P16, 2002; P17, 2005) modellieren ein
mehrstufiges Speicher-Modell mit einem Block-Transfer-Parameter, und Zhang (P27, ASPLOS 2025) bündelt
Prefetches hierarchisch über L1/L2/L3.

\subsection{Cache-Line-bewusstes Layout}\label{ssec:sota-cacheline}
Das cache-line-orientierte Layout (Abschnitt~\ref{ssec:cache-line-sizes}) ist der dominante
Leitgedanke der Layout-Theorie: Der CSS-Baum (P11, Rao/Ross 1999) führt den Cache-Line-Knoten ein, der
CSB\textsuperscript{+}-Baum (P12, Rao/Ross 2000) das zusammenhängende Geschwister-Cluster; Hankins und
Patel (P13, 2003) vermessen den Effekt der Knotengröße, Samuel et al.\ (P14, 2005) machen den Knoten
prozessor-bewusst. Der Adaptive Radix Tree (P01~\cite{leis2013art}) staffelt mit Node4/16/48/256 die
Knotengröße cache-line-orientiert, und Schmidt et al.\ (P32, 2025) optimieren für die HBM-Bandbreite.
Auf der Speicher-Layout-Achse stehen damit konkret: dichtes Array-of-Structs, cache-line-ausgerichtetes
AoS, spaltenorientiertes Struct-of-Arrays, bit-gepacktes LOUDS (P09, Jacobson 1989) und das
SIMD-gekachelte AoSoA.

\subsection{Adressübersetzung: TLB und dTLB}\label{ssec:sota-tlb}
Der dTLB (Abschnitt~\ref{ssec:tlb}) tritt im Stand der Technik vorwiegend als \emph{Mess-Ziel} auf ---
die Mess-Architektur führt eine eigene dTLB-Miss-Kategorie (Abschnitt~\ref{sec:measurement-basics}).
Algorithmenseitig adressieren ihn große Seiten (Huge-Pages), die mehrere Allokatoren (mimalloc,
jemalloc, TCMalloc, snmalloc) unterstützen, sowie Schmidt et al.\ (P32, 2025) empirisch mit
2-MiB-Seiten. Eine \emph{eigene} dTLB-Optimierungs-Achse bleibt in den Suchalgorithmus-Papern jedoch
unterbelegt --- eine Beobachtung, auf die die Forschungslücke (Abschnitt~\ref{sec:gap}) zurückkommt.

\subsection{Schreib-Protokolle, Update-Heuristiken und Inklusions-Politik}\label{ssec:sota-update}
Das in Abschnitt~\ref{ssec:coherence} eingeführte Cache-Update-Verhalten ist je ISA und je Cache-Ebene
konkret festgelegt. \emph{Schreib-Protokolle:} ARMv8 wählt Write-Through, Write-Back oder
Non-Cacheable sowie Read- bzw.\ Read-Write-Allocate über Seiten-Attribute~\cite{arm_arm}, während x86
als Normalfall Write-Back mit Write-Allocate verwendet~\cite{intel_sdm}. \emph{Heuristische
Ersetzungs-/Update-Strategien:} von LRU und Pseudo-LRU bis zur \emph{Re-Reference Interval Prediction}
(RRIP/DRRIP), die mit zwei Bit je Line scan- und thrash-resistente Verdrängung
erreicht~\cite{jaleel2010rrip} und auf der adaptiven Insertions-Politik aufbaut~\cite{qureshi2007adaptive}.
\emph{Inklusions-Politik:} Intel-Client-Prozessoren führen einen inklusiven L3, Intel-Server ab
Skylake-X einen nicht-inklusiven, während AMD Zen den L3 als überwiegend exklusiven Victim-Cache
betreibt~\cite{intel_optimization,amd_sog}. Diese drei Aspekte bestimmen gemeinsam, welche Ebenen ein
Schreibzugriff aktualisiert; da sie mikroarchitektonisch fixiert und zur Laufzeit nicht wählbar sind,
fließen sie als \emph{vermessene} Plattform-Eigenschaften in die systemoptimierte Binary ein
(Abschnitt~\ref{ssec:aware-oblivious}).

\subsection{Cache-bewusste und cache-oblivious Entwürfe}\label{ssec:sota-aware}
Beide Philosophien aus Abschnitt~\ref{ssec:aware-oblivious} sind im Stand der Technik instanziiert:
Bender et al.\ (P17, 2005) begründen den cache-oblivious B-Baum über das van-Emde-Boas-Layout, P16
(2002) erweitert das mehrstufige Layout-Modell, und Saikkonen/Soisalon-Soininen (P18, 2008; P19, 2016)
halten eine cache-sensitive Layout-Invariante auch unter Updates; Graefe/Larson (P15, 2001) vergleichen
beide Ansätze. Sämtliche praktischen Trie- und B\textsuperscript{+}-Indizes (P01--P14, P20) sind
hingegen \emph{cache-bewusst} ausgelegt. Die vorliegende Arbeit ordnet sich als dritter Weg ein
(Abschnitt~\ref{ssec:aware-oblivious}): gemessenes Lokalitätswissen statt fest verdrahteter oder
vollständig ignorierter Cache-Parameter.

% [AP-S1] ===== Ab hier ROHMATERIAL (vorlaeufige Nummerierung) — wird im naechsten Inkrement reorganisiert =====
% Die folgenden Cluster-Sektionen (A--F), die 3 Profil-/HW-Tabellen und die Forschungsluecke bleiben inhaltlich
% erhalten und werden in 3.3 (Achsen-Sezierung, je Achse Organ-Instanzen) / 3.4 (Workload-Frameworks) /
% 3.5 (Messen) / 3.6 (Architektur-Design) ueberfuehrt; die Forschungsluecke wird zu 3.7.

\section{Trie-Familie (Cluster A)}\label{sec:cluster-a}

\textbf{P01 ART}~\cite{leis2013art} führt den Adaptive Radix Tree mit adaptiven Knotengrößen
Node4/16/48/256 ein. \textbf{P02 HOT}~\cite{binna2018hot} optimiert die Trie-Höhe über
Multi-Byte-Präfixe und Bit-Vektor-Indizes. \textbf{P03 Masstree}~\cite{mao2012masstree}
kombiniert einen Trie mit B-Bäumen und einem Permutations-Index je interner Knoten.
\textbf{P04 CoCo-Trie}~\cite{boffa2024coco} verwendet succinct-kodierte kompakte Seiten.
\textbf{P05 START}~\cite{fent2020start} erweitert ART um Selbstoptimierung der
Knotenrepräsentation. P09 LOUDS (Jacobson 1989) ist die succinct-Trie-Grundlage, und
\textbf{P10 SuRF}~\cite{zhang2018surf} nutzt LOUDS zur Bereichsfilterung.

\section{Hybrid- und B\texorpdfstring{\textsuperscript{+}}{+}-Familie (Cluster B)}\label{sec:cluster-b}

\textbf{P06 B\textsuperscript{2}-Baum}~\cite{schmeisser2022b2tree} setzt einen zweistufigen
Verzweigungsgrad für kombinierte L1+L2-Cache-Bewusstheit ein. \textbf{P07
Wormhole}~\cite{wu2019wormhole} kombiniert eine Hash-Tabelle mit einer linearen Präfix-Sicht.
P11 CSS-Baum (Rao/Ross 1999) ist der erste Cache-Sensitive Search Tree; P12
CSB\textsuperscript{+}-Baum (Rao/Ross 2000) entwickelt ihn zu einem B\textsuperscript{+}-Baum
weiter; P13 Hankins/Patel (2003) untersucht den Effekt der Knotengröße auf die
Cache-Performance.

\section{Layout-Theorie (Cluster C)}\label{sec:cluster-c}

P14 Samuel/Pedersen/Bonnet (2005) macht den CSB\textsuperscript{+}-Baum prozessor-bewusst.
P15 Graefe/Larson (2001) untersucht B-Baum-Indizes gegenüber CPU-Caches. P16
Bender/Demaine/Farach-Colton (2002) führt cache-oblivious Baum-Layout (van Emde Boas) ein;
P17 Bender et al.\ (2005) erweitert es zu cache-oblivious B-Bäumen. P18
Saikkonen/Soisalon-Soininen (2008) und P19 (2016) entwickeln Multi-Level- und
layout-invariante B-Bäume.

\section{Prefetching (Cluster D und E)}\label{sec:cluster-de}

P20 ``B-Trees Are Back'' (Mueller/Benson/Leis 2025) demonstriert modernes B-Baum-Engineering.
P21 Chen/Gibbons/Mowry (2001) führt Prefetching-B\textsuperscript{+}-Bäume ein; P22 Chen et
al.\ (2002) entwickelt fraktales Prefetching. P23 Khan (2010) macht Cache-Prefetching
dynamisch adaptiv; P24 Naderan-Tahan/Sarbazi-Azad (2016) ergänzt einen adaptiven Filter auf
dem Cache-Prefetcher. P25 Mahling/Weisgut/Rabl (2025) kombiniert Bereichs-Scans mit
Hot-Path-Caching. P26 Zhang (FGCS 2024) und P27 Zhang (ASPLOS 2025, hierarchisches
Prefetching) sind die jüngsten Schemata, und P28 Kuehn (DaMoN 2023) motiviert datengetriebene
Cache-Optimierung.

\section{Synchronisation (Cluster F)}\label{sec:cluster-f}

P08 ``ART of Practical Synchronization'' (Leis 2016) führt das optimistische Lock-Coupling
(OLC) für ART ein. P29 RCU (McKenney, OLS 2001) und P30 Hazard Pointers (Michael 2004) sind
die kanonischen Reklamations-Schemata. P31 Ungethüm/Habich (2017) und P32 Schmidt/Habich
(2025) sind TU-Dresden-spezifische Hardware-Optimierungs-Überblicke; P33 (VAMPIR-Poster,
Berthold/Habich 2023) liefert die OTF2-Profiler-Integration.

\section{Allokator-Forschung (A01--A23)}\label{sec:allocators}

Die 21 Allokations-Paper sind systematisch in der Allokator-Matrix (Anhang~\ref{app:blocks})
katalogisiert. Rang-1-Repräsentanten sind Hoard (A01, ASPLOS 2000), Slab (A02, USENIX 1994),
Michael Lock-Free (A03, PLDI 2004), mimalloc (A04, MSR-TR-2019-18), jemalloc (A05), TCMalloc
(A06), snmalloc (A07, ISMM 2019) und NUMAlloc (A09, ISMM 2023). Rang-2/3 ergänzen CAMA (A12,
RTAS 2011), StarMalloc (A13, OOPSLA 2024), Crystalline-Reklamation (A17, PLDI 2024) und
PIM-malloc (A16, arXiv 2025).

\section{Profil-Schema mit dem \texorpdfstring{\texttt{<expected\_workload>}}{expected\_workload}-Tag}\label{sec:profile-schema}

Für jeden SOTA-Algorithmus und jeden Allokator hält ein Profil-XML in
\texttt{cache\_engine/algorithm\_profiles/} Metadaten fest, die Achsen-Belegung (page, node,
traversal, value\_handle, concurrency, allocator, prefetch, telemetry, isa, layout,
reclamation), eine Schlüssel-Wert-Signatur und einen optionalen
\texttt{<expected\_workload>}-Tag, der die \texttt{ExperimentDriver}-Heuristik
(Kapitel~\ref{ch:methodology}) überschreibt. Alle 30 SOTA-Profile und alle 10
Allokator-Profile sind getaggt.

\begin{table}[!htbp]
\caption{SOTA-Algorithmus-Profile mit Workload-Affinität}\label{tab:sota-profiles}
\centering\small
\begin{tabular}{lll}
\toprule
\textbf{Profil} & \textbf{Paper-Ref.} & \textbf{expected\_workload} \\
\midrule
art & P01 (Leis 2013) & YCSB\_C \\
hot & P02 (Binna 2018) & YCSB\_C \\
masstree & P03 (Mao 2012) & YCSB\_A \\
coco\_trie & P04 (Boffa 2024) & YCSB\_C \\
start & P05 (Fent 2020) & YCSB\_C \\
b2tree & P06 (Schmeisser 2022) & YCSB\_E \\
wormhole & P07 (Wu 2019) & YCSB\_A \\
surf & P10 (Zhang 2018) & YCSB\_A \\
css\_tree & P11 (Rao/Ross 1999) & YCSB\_C \\
csb\_tree & P12 (Rao/Ross 2000) & YCSB\_E \\
hankins & P13 (Hankins 2003) & YCSB\_C \\
samuel & P14 (Samuel 2005) & YCSB\_C \\
graefe & P15 (Graefe 2001) & YCSB\_C \\
bender\_treelayout & P16 (Bender 2002) & YCSB\_C \\
bender\_cacheoblivious & P17 (Bender 2005) & YCSB\_C \\
saikkonen\_multilevel & P18 (Saikkonen 2008) & YCSB\_E \\
saikkonen\_layoutinvariant & P19 (Saikkonen 2016) & YCSB\_E \\
btreesareback & P20 (Mueller 2025) & YCSB\_E \\
chen\_prefetch & P21 (Chen 2001) & YCSB\_A \\
chen\_fractal & P22 (Chen 2002) & YCSB\_A \\
khan\_adaptive & P23 (Khan 2010) & YCSB\_A \\
naderan\_tahan & P24 (Naderan-Tahan 2016) & YCSB\_A \\
mahling & P25 (Mahling 2025) & YCSB\_E \\
zhang\_fgcs & P26 (Zhang 2024) & YCSB\_A \\
zhang\_asplos & P27 (Zhang 2025) & YCSB\_A \\
kuehn & P28 (Kuehn 2023) & YCSB\_C \\
rcu & P29 (McKenney 2001) & YCSB\_B \\
hazard\_pointers & P30 (Michael 2004) & YCSB\_B \\
ungethuem & P31 (Ungethuem 2017) & YCSB\_C \\
to\_stride & P32 (Schmidt 2025) & YCSB\_C \\
\bottomrule
\end{tabular}
\end{table}

Die Workloads verteilen sich auf 13$\times$ YCSB\_C (leselastig), 9$\times$ YCSB\_A
(Zipf-verteiltes Lesen+Aktualisieren), 6$\times$ YCSB\_E (Bereichs-Scans) und 2$\times$
YCSB\_B (95/5 Lesen/Aktualisieren). PRT-ART selbst ist mit YCSB\_F (Read-Modify-Write, für
DensityTracker-Aktualisierungen) getaggt. Tabelle~\ref{tab:allocator-profiles} listet die
zehn implementierten Allokator-Profile.

\begin{table}[!htbp]
\caption{Allokator-Profile mit Lizenz und Workload-Affinität}\label{tab:allocator-profiles}
\centering\small
\begin{tabular}{llll}
\toprule
\textbf{Profil} & \textbf{Familie} & \textbf{Lizenz} & \textbf{expected\_workload} \\
\midrule
hoard & A01 & Apache-2.0 & YCSB\_A \\
michael\_lockfree & A03 & BSD-3 & YCSB\_B \\
mimalloc & A04 & MIT & YCSB\_A \\
jemalloc & A05 & BSD-2 & YCSB\_B \\
tcmalloc & A06 & BSD-3 & YCSB\_C \\
snmalloc & A07 & MIT & YCSB\_A \\
scalloc & A08 & BSD-3 & YCSB\_A \\
rpmalloc & A10 & Public Domain & YCSB\_C \\
lrmalloc & A11 & BSD-3 & YCSB\_B \\
dlmalloc & A20 & CC0 & YCSB\_C \\
\bottomrule
\end{tabular}
\end{table}

\section{Hardware- und Scheduling-Strategie pro Paper}\label{sec:hw-sched}

Nach der N-Phasen-Erweiterung der Bausteine-Matrix (Abschnitt~\ref{sec:axes}) sind zwei
zusätzliche Algorithmus-Permutations-Achsen verankert:

\begin{itemize}
  \item \textbf{Achse~12 Hardware-Strategie}: welche Hardware-Merkmale der Algorithmus aktiv
        nutzt --- Sub-Achsen 12.1 SIMD-Familie (AVX2/AVX-512/NEON/SVE2/scalar), 12.2
        Cache-Level-Targeting (L1/L2/L3/HBM-aware), 12.3 NUMA-Strategie, 12.4
        Prefetch-Hardware (PREFETCH/PREFETCHNTA/PREFETCHW), 12.5 Atomic-Instruction-Familie
        (CAS/LL-SC/RMW-extended).
  \item \textbf{Achse~13 Scheduling-Strategie}: wie Arbeit verteilt wird --- Sub-Achsen 13.1
        Worker-Pool-Layout (thread-per-core / work-stealing / CPU-pinning), 13.2
        SIMD-Worker-Anzahl-Limit (hardwarebedingt typisch 2 von $N$ Kernen), 13.3
        Heterogeneous-Core-Dispatch (Intel-Hybrid P-/E-Kerne), 13.4 Co-Routine-Strategie,
        13.5 Batch-Granularität.
\end{itemize}

Tabelle~\ref{tab:hw-sched} zeigt eine Auswahl dieser zwei Achsen je SOTA-Paper; der
Vollkatalog steht in Anhang~\ref{app:blocks}.

\begin{table}[h!]
\centering\small
\caption{Hardware- und Scheduling-Strategie pro SOTA-Paper (Auswahl)}\label{tab:hw-sched}
\begin{tabular}{lll}
\toprule
\textbf{Paper} & \textbf{Hardware-Strategie (Achse 12)} & \textbf{Scheduling-Strategie (Achse 13)} \\
\midrule
P01 ART & AVX2, L1-aware, none, CAS & thread-per-core \\
P02 HOT & AVX2 (Bit-Vektor), L1-aware, CAS & thread-per-core \\
P03 Masstree & scalar, L1-aware, CAS & thread-per-core + work-stealing \\
P05 START & AVX2 (Multi-Byte-Span), L1-aware, CAS & thread-per-core \\
P20 LeanStore & AVX2, HBM-aware, NUMA-aware, CAS & work-stealing + hybrid \\
P27 hp-soft & PREFETCH (L1/L2/L3-Bundle), all-cache, CAS & thread-per-core + Co-Routine \\
P28 Kuehn & PREFETCH (scalar Counter), L1-aware, CAS & thread-per-core + Sampling \\
P31 Ungethuem & AVX-512 (HBM), HBM-aware, NUMA-aware, CAS & hybrid (P-/E-Kerne) \\
P32 Schmidt & AVX-512 (SIMD-Stride), HBM-aware, NUMA-aware, CAS & hybrid \\
\midrule
\textbf{PRT-ART} & \textbf{AVX2, L1-aware (TLB-aligned), local NUMA, PREFETCH, CAS (OLC)} & \textbf{thread-per-core + 2-SIMD-Worker-Limit + hybrid + Co-Routine + micro-batch} \\
\bottomrule
\end{tabular}
\end{table}

\section{Forschungslücke}\label{sec:gap}

Jede der untersuchten Arbeiten optimiert eine oder wenige Achsen der Cache-Hierarchie
isoliert, und die Ergebnisse werden in nicht vergleichbaren Settings berichtet. Was fehlt ---
und was diese Arbeit beiträgt --- ist eine systematische, orthogonale \emph{Achsen-Zerlegung}
dieser Entwürfe zusammen mit einer permutierten Mess-Architektur, die sie als explizite
Konfigurationen rekonstruiert und auf einer gemeinsamen Basis vermisst.
